%% ============================================================
\section{Introduction}\label{sec:intro}
%% ============================================================

%% ============================================================
%% Figure 1: main overview figure
%% ============================================================
\begin{figure}[t]
\centering
\begin{tikzpicture}[
  font=\scriptsize,
  box/.style={draw, rounded corners=2pt, align=center, inner sep=3pt},
  stage/.style={box, fill=codeblue!8, draw=codeblue!60, thick},
  neg/.style={box, fill=codered!6, draw=codered!50},
  chip/.style={box, fill=white, draw=black!40, inner sep=2pt, font=\tiny},
  arr/.style={-{Stealth[length=2mm]}, thick, black!70},
  lbl/.style={font=\tiny, black!60},
  >=stealth,
]
%% ---------------- panel (a): task ----------------
\node[anchor=north west] (atitle) at (0,5.6) {\textbf{(a) Closed-book task}};
\node[box, anchor=north west, fill=white, font=\ttfamily\tiny, align=left]
  (prog) at (0,5.1) {
    requires n >= 0;\\
    x = 0; y = n;\\
    while (y > 0)\\
    \{ x++; y--; \}\\
    \textcolor{codered}{\sout{assert x == n;}}
  };
\node[lbl, anchor=west] at ($(prog.south west)+(0,-0.22)$) {postcondition \textcolor{codered}{stripped}};
\node[stage, anchor=north] (pol) at ($(prog.south)+(0,-0.65)$) {policy $\pi_\theta$};
\node[box, anchor=north, fill=codegreen!8, draw=codegreen!60, font=\ttfamily\tiny]
  (summ) at ($(pol.south)+(0,-0.45)$) {x + y == n;\\ 0 <= y <= n;};
\node[lbl, anchor=north] at ($(summ.south)+(0,-0.05)$) {inductive loop summary};
\draw[arr] (prog) -- (pol);
\draw[arr] (pol) -- (summ);

%% ---------------- panel (b): sampler ----------------
\node[anchor=north west] (btitle) at (4.9,5.6) {\textbf{(b) Sampler: positives + truthful negatives}};
\begin{scope}[shift={(5.7,1.55)}]
  % axes
  \draw[->, black!50] (0,0) -- (4.35,0) node[right, lbl] {$x$};
  \draw[->, black!50] (0,0) -- (0,3.0) node[above, lbl] {$y$};
  % observed range box
  \draw[dashed, black!35] (-0.18,-0.18) rectangle (2.55,2.68);
  \node[lbl, anchor=west] at (1.45,2.5) {observed range};
  % manifold x+y=n, extended past the exit
  \draw[codeblue!40, dashed] (0,2.45) -- (3.62,-1.17);
  % positives on the manifold
  \foreach \i in {0,...,6} {
    \fill[codeblue] ({0.35*\i},{2.45-0.35*\i}) circle (1.6pt);
  }
  \node[lbl, text=codeblue!80!black, anchor=west] at (0.35,2.9) {positives: $x{+}y{=}n$, $0 \le y \le n$};
  % relation negative: off the manifold
  \node[text=codered] at (0.75,1.05) {$\times$};
  \node[lbl, text=codered] at (0.75,0.72) {relation $\pm\{1,2\}$};
  \node[lbl, text=codered] at (0.72,0.47) {off the law};
  % escape: push one variable out of range, dotted displacement arrow
  \draw[densely dotted, codered!70, -{Stealth[length=1.4mm]}] (1.15,1.4) -- (3.75,1.4);
  \node[text=codered] at (3.87,1.4) {$\blacktriangle$};
  \node[lbl, text=codered, anchor=east] at (4.25,1.72) {escape: out of range,};
  \node[lbl, text=codered, anchor=east] at (4.25,1.02) {far, seed-hashed};
  % over-run continuation: on the manifold, past the exit
  \foreach \i in {8,9} {
    \draw[codered, thick] ({0.35*\i-0.06},{2.45-0.35*\i-0.06}) rectangle ({0.35*\i+0.06},{2.45-0.35*\i+0.06});
  }
  \node[lbl, text=codered, anchor=west] at (0.0,-0.72) {over-run: past the exit};
\end{scope}
% filter strip
\node[neg, anchor=north west] (f1) at (5.35,0.28) {reachable?};
\node[neg, anchor=west] (f2) at ($(f1.east)+(0.18,0)$) {entry-feasible?};
\node[neg, anchor=west] (f3) at ($(f2.east)+(0.18,0)$) {tainted?};
\node[lbl, anchor=west] at ($(f3.east)+(0.15,0)$) {$\Rightarrow$ veto};
\node[lbl, anchor=north west] at (5.35,-0.25) {trace units: one impossible history $=$ one unit};

%% ---------------- panel (c): reward ----------------
\node[anchor=north west] (ctitle) at (10.9,5.6) {\textbf{(c) Hack-resistant dense reward}};
\node[box, anchor=north west, font=\ttfamily\tiny, align=left] (roll) at (10.9,5.1)
  {rollout $A$: clause set};
\node[stage, anchor=north west, align=left] (casc) at ($(roll.south west)+(0,-0.3)$)
  {complexity gate $\to$ scope\\ $\to$ positive soundness\\ $\to$ Houdini induction};
\node[lbl, anchor=west] at ($(casc.east)+(0.1,0.25)$) {$\mathrm{surv}(A)$};
\node[box, anchor=north west, fill=white] (score) at ($(casc.south west)+(0,-0.35)$)
  {$r = \max(0,\, \tfrac12\,\mathrm{base} + \tfrac12\,\mathrm{marg} - \tfrac1{20}\,\mathrm{junk})$};
\node[chip, anchor=north west] (s0) at ($(score.south west)+(0.05,-0.3)$) {seed 0};
\node[chip, anchor=west] (s1) at ($(s0.east)+(0.12,0)$) {seed 1};
\node[chip, anchor=west, draw=codered!60, fill=codered!6] (sh) at ($(s1.east)+(0.12,0)$) {holdout};
\node[box, anchor=west, fill=codegreen!8, draw=codegreen!60] (min) at ($(sh.east)+(0.25,0)$) {$\min$};
\draw[arr] (roll) -- (casc);
\draw[arr] (casc) -- (score);
\draw[arr] ($(score.south)+(0,-0.05)$) -- ($(score.south)+(0,-0.28)$);
\draw[arr] (min.east) -- ++(0.35,0) node[anchor=west, lbl, align=left] {GRPO\\ update};
\end{tikzpicture}
\caption{\textbf{\sam{} overview.}
(a)~The policy sees only precondition, guard, and body---the postcondition is stripped---and must emit the \emph{strongest} inductive loop summary, not an invariant tailored to a given target.
(b)~The sampler executes the instrumented loop over a deterministic, precondition-respecting input sweep; positives are observed loop-head states, and negatives are \emph{truthful impossible traces} in two dimensions: relational perturbations off densely sampled trace windows, and bound violations (over-run continuations past a genuine exit, and seed-hashed range escapes).
Three vetoes (reachability, entry feasibility, nondeterminism taint) guarantee every emitted negative is unreachable under its input.
(c)~Each rollout is reduced to its surviving inductive core, scored by the fraction of impossible traces it rejects (plus ablation-based marginal credit and a junk discount), and \emph{minimized} across canonical and held-out sampling seeds, so memorizing any finite sample earns nothing.}
\label{fig:overview}
\end{figure}


Loop invariants are the linchpin of deductive program verification.
Ever since the inductive-assertion method of Floyd~\citep{floyd1967assigning} and the axiomatic semantics of Hoare~\citep{hoare1969axiomatic}, verifying a loop $\mathtt{while}\,(B)\,S$ against a contract $\{P\}\cdot\{Q\}$ has reduced to exhibiting an \emph{inductive invariant} $I$ satisfying initiation ($P \Rightarrow I$), consecution ($\{I \wedge B\}\,S\,\{I\}$), and safety ($I \wedge \neg B \Rightarrow Q$).
Finding $I$ automatically has consequently been one of the most intensively studied problems in formal methods.
Four decades of research have attacked it with abstract interpretation~\citep{cousot1977abstract,cousot1978automatic,karr1976affine}, constraint- and template-based synthesis~\citep{colon2003linear,sankaranarayanan2004non,gulwani2008program,gupta2009invgen}, interpolation and IC3/PDR-style generalization~\citep{mcmillan2003interpolation,bradley2011sat,komuravelli2014smt}, abductive inference~\citep{dillig2013inductive}, syntax-guided and stochastic search~\citep{alur2013sygus,fedyukovich2018syntax}, data-driven learning from program states~\citep{sharma2012interpolants,sharma2013data,garg2014ice,garg2016learning,padhi2016pie,zhu2018data,ezudheen2018horn}, neural architectures~\citep{si2018learning,si2020code2inv,ryan2020cln2inv,yao2020gcln}, and, most recently, large language models~\citep{pei2023can,kamath2023finding,chakraborty2023ranking,wu2024lemur,wu2024lam4inv,liu2024acinv,wen2024enchanting,liu2024towards}.

\paragraph{The target-dependence problem.}
Despite their diversity, these lines of work share a common task definition: the input is a program \emph{together with a verification target}---a postcondition, assertion, or safety property $Q$---and the output is an invariant just strong enough to discharge $Q$.
The target is not incidental; it is load-bearing.
CEGIS- and ICE-style learners derive their counterexamples from failed attempts to prove $Q$~\citep{garg2014ice,garg2016learning,padhi2016pie}; Code2Inv shapes its reinforcement signal from the solver's rejection of $Q$-directed verification conditions~\citep{si2018learning,si2020code2inv}; LLM pipelines put the assertion in the prompt and score candidates by whether the verifier closes the goal~\citep{kamath2023finding,chakraborty2023ranking,wu2024lemur,wu2024lam4inv,liu2024acinv}.

We contend that this framing rests on an uncomfortable circularity.
First, verification targets are \emph{scarce}: real code overwhelmingly ships without formal postconditions, and writing them is itself the specification bottleneck that specification mining and, lately, LLM-based specification generation have tried to relieve~\citep{daikon,endres2024can,ma2024specgen,wen2024enchanting}.
Second, and more fundamentally, for loop programs the verification target is very often nothing more than \emph{the invariant itself evaluated at loop exit}.
The safety condition $I \wedge \neg B \Rightarrow Q$ means that the natural $Q$ for a loop---the one a human or a benchmark author writes down---is the exit-time projection of an invariant the author already had in mind; Furia and Meyer make the reverse direction explicit by \emph{deriving} candidate invariants syntactically from postconditions~\citep{furia2010inferring}.
A system that requires $Q$ in order to find $I$ therefore presupposes the very reasoning it is supposed to automate.
Target-directed training compounds the problem: when the reward is ``did the verifier close $Q$,'' the learner is incentivized to produce the \emph{weakest} invariant sufficient for the given target---frequently a light rewriting of $Q$ itself---which transfers to no other property of the same loop.

\paragraph{Our reframing: inductive loop summaries.}
We instead study the task of generating an inductive \emph{loop summary}: given only the precondition, guard, and body---\emph{no postcondition}---produce the strongest practically expressible inductive invariant characterizing the loop's reachable states at the loop head (Figure~\ref{fig:overview}a).
This framing follows the tradition of loop summarization in static analysis and symbolic execution~\citep{kroening2008loop,godefroid2011automatic,xie2016proteus,farzan2015compositional,kincaid2017compositional}, but asks for summaries in the standard inductive-invariant form consumed by deductive verifiers (ACSL, checked by Frama-C/WP~\citep{framac,acsl}).
A tight summary is strictly more useful than a target-directed invariant: it discharges \emph{every} downstream assertion $Q$ with $I \wedge \neg B \Rightarrow Q$, it is reusable across changing specifications, and generating it exercises exactly the semantic reasoning---conservation laws, relational couplings, progress bounds---that the target-directed setting lets a system shortcut.

\paragraph{The reward problem, and why it is hard.}
Dropping the target, however, destroys the training signal every prior learning-based system relied on: with no $Q$ there is no binary verified/failed outcome to optimize.
``Is $I$ inductive?'' alone is useless as a reward---$\mathit{true}$, the guard, and any tautology are inductive.
Recent work concurs that the verifier bit is too coarse a signal for learning even in the target-directed setting: InvBench scores invariants by wall-clock speedup rather than by verification alone~\citep{wei2025invbench}, and \citet{pinto2026notall} show that grading and curating invariants by quality, not mere validity, is what makes them useful training data.
What must be rewarded is \emph{soundness plus empirical tightness}: how effectively does a proved invariant separate sampled reachable states from audited synthetic countertraces?
We operationalize this with an execution-grounded dense reward (Figure~\ref{fig:overview}b,c).
A sampler compiles and runs the instrumented loop over a deterministic input sweep to collect \emph{positive} loop-head states, then constructs three families of synthetic negative candidates: relational perturbations around densely observed trace neighborhoods, continuations of the real body past a genuine exit, and ladder-step escapes beyond a variable's sampled range.
A candidate rollout is first reduced to its inductive core by Houdini-style pruning~\citep{flanagan2001houdini}, with Frama-C/WP as the authoritative backend, and the surviving subset is scored by the fraction of negative traces it rejects, plus a marginal (ablation) term for group-relative credit assignment.

Like any finite-sample signal, such a reward is vulnerable to mislabeled examples and sample-specific predicates~\citep{amodei2016concrete,skalse2022defining,gao2023scaling}.
We therefore keep the trusted mechanism small.  Construction-time vetoes discard a candidate negative if it has been observed reachable or could be a fresh loop entry under admissible inputs, and oracle-controlled loops emit no synthetic negatives; real Houdini then removes candidate invariants that merely fit the finite positive sample but are not inductive.  The scoring rule itself has only two terms: $0.5$ times standalone rejection and $0.5$ times group-ablation marginal rejection.

\paragraph{Contributions.}
This paper makes the following contributions:
\begin{itemize}[leftmargin=1.5em, itemsep=2pt, topsep=2pt]
  \item \textbf{Task reframing.} We identify the target-dependence circularity in loop invariant generation and formulate the alternative task of \emph{closed-book inductive loop summary generation}, which subsumes target-directed generation ($\S$\ref{sec:overview}).
  \item \textbf{An execution-grounded, minimal reward.} We design a dense reward from three audited negative-candidate families and conservative construction guards, with Houdini-pruned survival, trace-unit accounting, and the complete formula $r=0.5\,\mathrm{base}+0.5\,\mathrm{marginal}$ ($\S$\ref{sec:method}).
  \item \textbf{An open three-component pipeline.} \sam{} decomposes into an example sampler, a self-contained HTTP reward service whose image bundles Frama-C/WP, and an independent inference framework that performs no positive/negative sampling and optionally runs verdict-driven multi-round refinement before final certification ($\S$\ref{sec:overview}).
  \item \textbf{Direct validation of discrimination and labels.} The maintained discrimination harness ranks known-quality rollout families on 10 deterministic programs, and the full mislabel audit found zero hard collisions among 486,302 negative witnesses from all 366 benchmark programs ($\S$\ref{sec:experiments}).
\end{itemize}
